\chapter{The Derivative: Slope, Rate, and the Tangent}

In Chapter~9, we stepped far back to study the end behavior of functions as $x\to\pm\infty$. Now we reverse the lens and zoom all the way in on a single point, returning to the question we have been circling since Chapter~1:

\begin{focusbox}
How fast is something changing \emph{right now} -- at a single instant, not over
a stretch of time? And what does it mean for a quantity to have a rate
of change at one instant, when ``rate'' seems to require two moments to compare?
\end{focusbox}

The answer is the \textbf{derivative} -- a single number measuring how
sensitively a function's output reacts to a tiny nudge in its input. It is, all
at once, three things that turn out to be the same thing: the slope of a tangent
line, an instantaneous rate of change, and the best straight-line approximation
to a curve. By the end of this chapter you will see why those three are one,
you will compute a few derivatives from raw first principles with nothing but
algebra and a limit, and you will understand exactly what the famous symbol
$\dv{y}{x}$ is really saying -- and what it is \emph{not} saying.

\chapterline

\section{The problem of the instant}
\label{ch10:problem-of-instant}

Let us start with the most physical version of the question, the one that drove
Newton. Drop a ball. In the absence of air resistance, the distance it has fallen
after $t$ seconds is, to good accuracy,
\[
  s(t) = 16t^2 \quad \text{(feet, with $t$ in seconds).}
\]
After $1$ second it has fallen $16$ feet; after $2$ seconds, $64$ feet; after
$3$ seconds, $144$ feet. Now I ask: \emph{how fast is the ball moving at the
instant $t = 2$?}

You know how to find an \emph{average} speed. Average speed over a time interval
is distance covered divided by time elapsed. From $t = 2$ to $t = 3$ the ball
falls $144 - 64 = 80$ feet in $1$ second, so its average speed over that second
is $80$ feet per second. But that is the average over a whole second, during
which the ball was speeding up the entire time. At the \emph{start} of that second
it was slower than $80$; at the \emph{end}, faster. The number $80$ is not the
speed ``at $t = 2$.'' It is a blur of an entire second's worth of speeds.

\begin{intuitionbox}
To get the speed \emph{at} an instant, shrink the interval. Don't measure over a
whole second -- measure over a tenth of a second, then a hundredth, then a
millionth. The shorter the interval, the less the speed has a chance to change
across it, so the average over that tiny interval gets closer and closer to the
``true'' instantaneous speed. The instantaneous speed is what those averages are
\emph{heading toward}. That is a limit.
\end{intuitionbox}

Let us actually watch them head somewhere. Take the interval from $t = 2$ to
$t = 2 + h$, where $h$ is a small time step. The distance fallen during that
interval is $s(2 + h) - s(2)$, and the average speed is
\[
  \frac{s(2 + h) - s(2)}{h}
  = \frac{16(2 + h)^2 - 16(2)^2}{h}.
\]
Expand the top:
\[
  16(2 + h)^2 - 16(4)
  = 16\bigl(4 + 4h + h^2\bigr) - 64
  = 64 + 64h + 16h^2 - 64
  = 64h + 16h^2.
\]
Divide by $h$ (we can cancel because $h$, while
small, is not yet zero):
\[
  \frac{64h + 16h^2}{h} = 64 + 16h.
\]
So the average speed over the interval $[2,\, 2+h]$ is exactly $64 + 16h$ feet
per second. Watch what happens as $h$ shrinks.

\begin{center}
\begin{tabular}{cc}
\toprule
$h$ (seconds) & average speed $64 + 16h$ (ft/s) \\
\midrule
$1$        & $80$ \\
$0.1$      & $65.6$ \\
$0.01$     & $64.16$ \\
$0.001$    & $64.016$ \\
$0.0001$   & $64.0016$ \\
\bottomrule
\end{tabular}
\end{center}

The averages are marching toward $64$. And we don't have to guess -- the algebra
told us the average is $64 + 16h$, so as $h \to 0$ the $16h$ term vanishes and
the limit is exactly $64$. \emph{The ball is moving $64$ feet per second at the
instant $t = 2$.} We never divided by zero. We divided by a small nonzero $h$,
simplified, and only \emph{then} let $h$ go to zero.

\begin{walkbox}
\textbf{Pause and think:} What happens if you measure the average speed by
stepping \emph{backward} in time instead of forward, over the interval $[2-h, 2]$?
Take $h = 0.01$, so the interval is $[1.99, 2]$. Compute:
\[
  \frac{s(2) - s(1.99)}{0.01} = \frac{64 - 16(1.99)^2}{0.01}
  = \frac{64 - 63.3616}{0.01} = \frac{0.6384}{0.01} = 63.84\text{ ft/s}.
\]
Notice that $63.84 = 64 - 16(0.01)$. Stepping forward from $2$ gave $64 + 16h$ (speeds
just above $64$); stepping backward gives $64 - 16h$ (speeds just below $64$). As
$h \to 0$ from either direction, both squeeze onto the identical value: $64\text{ ft/s}$.
\end{walkbox}

\begin{happeningbox}
Notice what the limit accomplished. The raw fraction
$\dfrac{s(2+h) - s(2)}{h}$ is the indeterminate form $\frac{0}{0}$ when $h = 0$:
the ball travels zero distance in zero time, and $\frac{0}{0}$ is meaningless. The
instantaneous speed is \emph{not} ``$\frac{0}{0}$ evaluated.'' It is the
\emph{number that the fraction approaches} as $h$ shrinks -- a number the fraction
itself never literally equals at $h=0$, but circles arbitrarily closely. The limit
is precisely the tool that extracts a definite value out of a $\frac{0}{0}$ that
would otherwise be nonsense. Without limits, the instant has no speed. With limits,
it has exactly one.
\end{happeningbox}

\section{The same picture, drawn: secants closing in on a tangent}
\label{ch10:secant-tangent}

That falling-ball calculation has a geometric twin, and once you see the picture
you will never lose it. Plot the function -- any function $y = f(x)$ -- as a curve.
Pick a point $P = (a, f(a))$ on it. We want the steepness of the curve \emph{at}
$P$: the slope of the line that just grazes the curve there, touching it at $P$
and matching its direction. That grazing line is the \textbf{tangent line} -- a
straight line that touches the curve at that point without cutting across it
locally, matching the curve's exact direction at the contact point.

The trouble is that slope needs \emph{two} points -- rise over run -- and the
tangent touches at only one. So we approximate using two points, exactly as before. Pick a second,
nearby point $Q = (a + h,\, f(a + h))$ on the curve, and draw the line through
$P$ and $Q$. A line through two points of a curve is called a \textbf{secant}
line -- a straight chord connecting two separate points on the graph -- and
its slope we can compute directly from two points:
\[
  \text{slope of secant } PQ
  = \frac{\text{rise}}{\text{run}}
  = \frac{f(a + h) - f(a)}{(a + h) - a}
  = \frac{f(a + h) - f(a)}{h}.
\]
This fraction is so important it has a name: the \textbf{difference quotient} --
the ratio comparing the vertical change in output to the horizontal step in input
over a finite interval. It is the average rate of change of $f$ over the step
from $a$ to $a + h$ -- the geometric face of the very same average speed we
computed for the ball. Figure~\ref{ch10:fig-diff-quotient} lays out this
rise-over-run triangle directly on a curve.

\begin{figure}[htbp]
\centering
\begin{tikzpicture}
\begin{axis}[
  width=11.5cm, height=7cm,
  title style={font=\small},
  title={The geometry of the difference quotient: rise over run on a secant chord},
  axis lines=middle,
  xlabel={$x$}, ylabel={$y$},
  xmin=-0.3, xmax=3.4, ymin=-0.4, ymax=4.2,
  xtick={1, 2.5}, xticklabels={$a$, $a+h$},
  ytick={0.8, 2.9}, yticklabels={$f(a)$, $f(a+h)$},
  domain=0.2:3.1, samples=80,
  clip=false,
]
\addplot[curveblue, very thick] {0.4 + 0.4*x^2};
\node[curveblue, above left] at (axis cs:3.0, 3.8) {$y = f(x)$};

\coordinate (P) at (axis cs:1, 0.8);
\coordinate (Q) at (axis cs:2.5, 2.9);
\coordinate (R) at (axis cs:2.5, 0.8);

\addplot[gray, dotted, thick] coordinates {(1,0) (1,0.8) (0,0.8)};
\addplot[gray, dotted, thick] coordinates {(2.5,0) (2.5,2.9) (0,2.9)};

\draw[quotientorange, thick, dashed] (P) -- (R);
\draw[quotientorange, thick, dashed] (R) -- (Q);

\addplot[quotientorange, very thick, domain=0.4:3.0] {0.8 + 1.4*(x - 1)};

\addplot[only marks, mark=*, mark size=2pt, black] coordinates {(1, 0.8)};
\node[above left] at (axis cs:1, 0.8) {$P\bigl(a, f(a)\bigr)$};
\addplot[only marks, mark=*, mark size=2pt, quotientorange] coordinates {(2.5, 2.9)};
\node[above left] at (axis cs:2.5, 2.9) {$Q\bigl(a+h, f(a+h)\bigr)$};

\node[below, quotientorange, font=\small] at (axis cs:1.75, 0.8) {$\text{run} = h$};
\node[right, quotientorange, font=\small] at (axis cs:2.5, 1.85) {$\text{rise} = f(a+h) - f(a)$};

\node[quotientorange!80!black, above, sloped, font=\small] at (axis cs:1.75, 1.85) {$\text{secant line}$};
\end{axis}
\end{tikzpicture}
\caption{The geometric anatomy of the difference quotient. The secant line connects
$P=\bigl(a, f(a)\bigr)$ and $Q=\bigl(a+h, f(a+h)\bigr)$. Its slope is the vertical rise
$f(a+h)-f(a)$ divided by the horizontal run $h$. When $h$ shrinks toward zero, $Q$
slides along the curve toward $P$, and this secant line pivots toward the tangent.}
\label{ch10:fig-diff-quotient}
\end{figure}

Now slide $Q$ along the curve toward $P$. As $Q$ approaches $P$, the secant line
pivots about $P$, and its slope changes. The picture practically begs you to
believe that the secants settle down toward one limiting line -- the tangent --
and that the secant slopes settle toward one limiting number -- the slope of the
tangent. Figure~\ref{ch10:fig-secants} shows three secants for shrinking $h$,
each tilting closer to the dashed tangent.

\begin{figure}[htbp]
\centering
\begin{tikzpicture}
\begin{axis}[
  width=12cm, height=8cm,
  title style={font=\small},
  title={Secant-to-tangent convergence: chords pivoting as $h \to 0$},
  axis lines=middle,
  xlabel={$x$}, ylabel={$y$},
  xmin=-0.3, xmax=3.2, ymin=-0.5, ymax=6,
  xtick={1,2,3}, ytick={\empty},
  domain=0:3, samples=100,
  clip=false,
]
% the curve y = x^2/2
\addplot[curveblue, very thick] {x^2/2};
% point P at a = 1
\addplot[only marks, mark=*, mark size=2pt, black] coordinates {(1,0.5)};
\node[below right] at (axis cs:1,0.5) {$P$};
% secant 1: Q at x=3 -> slope (4.5-0.5)/2 = 2
\addplot[quotientorange, thick, domain=-0.2:3.2] {0.5 + 2*(x-1)};
\addplot[only marks, mark=*, mark size=1.5pt, quotientorange] coordinates {(3,4.5)};
\node[above left] at (axis cs:3,4.5) {$Q_1$};
% secant 2: Q at x=2 -> slope (2-0.5)/1 = 1.5
\addplot[quotientorange!70, thick, domain=-0.2:3] {0.5 + 1.5*(x-1)};
\addplot[only marks, mark=*, mark size=1.5pt, quotientorange] coordinates {(2,2)};
\node[above left, fill=white, inner sep=1pt] at (axis cs:2,2) {$Q_2$};
% secant 3: Q at x=1.5 -> slope (1.125-0.5)/0.5 = 1.25
\addplot[quotientorange!45, thick, domain=-0.2:2.5] {0.5 + 1.25*(x-1)};
\addplot[only marks, mark=*, mark size=1.5pt, quotientorange] coordinates {(1.5,1.125)};
\node[quotientorange!90!black, font=\small] (lblQ3) at (axis cs:1.2, 2.0) {$Q_3$};
\draw[->, quotientorange!80!black, shorten >=2pt] (lblQ3) -- (axis cs:1.48, 1.18);
% tangent: slope 1 at P
\addplot[tangentred, very thick, dashed, domain=-0.2:3] {0.5 + 1*(x-1)};
\end{axis}
\end{tikzpicture}
\caption{Secant lines through $P=(1,\tfrac12)$ and points $Q_1,Q_2,Q_3$ on
$y=\tfrac12 x^2$, with $Q$ sliding toward $P$. The secant slopes ($2$, then
$1.5$, then $1.25$) close in on the slope of the dashed tangent line, which is
$1$. The tangent is the limit of the secants.}
\label{ch10:fig-secants}
\end{figure}

\begin{walkbox}
Before reading on, picture the secant pivoting and ask yourself: \emph{why}
should the secants approach a single line at all? What could go wrong? Hold a
ruler against a smooth bowl and tilt it as the second contact point slides in --
you can feel it lock onto one direction. Now imagine the curve has a sharp corner
at $P$ (think of the letter ``V''). Slide $Q$ in from the right and you get one
limiting slope; slide it in from the left and you get a different one. The secants
do \emph{not} agree on a single line. We will meet this failure again -- it is
exactly where derivatives stop existing.
\end{walkbox}

The two stories -- falling ball, grazing line -- are not mere analogies.
They are mathematically identical calculations. ``Instantaneous speed'' is the
limit of average speeds; ``slope of tangent'' is the limit of secant slopes; and
average speed on a position-time graph is the slope of a secant line.

\section{The definition of the derivative}
\label{ch10:def}

We have now done the same thing twice, so let us crystallize it once and name it.

\begin{definition}[The derivative at a point]
\label{ch10:def-derivative}
Let $f$ be a function defined on an open interval containing the point $a$
(meaning $f$ is defined on both sides of $a$, with a small buffer of space around
it). The \textbf{derivative of $f$ at $a$}, written $f'(a)$, is
\[
  f'(a) \;=\; \lim_{h \to 0} \frac{f(a + h) - f(a)}{h},
\]
\emph{provided this limit exists}. When it does, we say $f$ is
\textbf{differentiable at $a$}. The number $f'(a)$ is the slope of the tangent
line to the graph $y = f(x)$ at the point $(a, f(a))$, and equivalently the
instantaneous rate of change of $f$ with respect to $x$ at $x = a$.
\end{definition}

Two small but important remarks before we use it.

\begin{remark}
The phrase ``provided this limit exists'' is doing real work. The derivative is a
\emph{limit}, and limits can fail to exist (Chapters 5--7). When the limit does
not exist, $f$ has no derivative at $a$ -- no well-defined slope, no instantaneous
rate. That is not a defect in our method; it is a true fact about the function,
and Section~\ref{ch10:diff-vs-cont} is devoted to what such failures look like.
(Why must $f$ be defined on an open interval around $a$? Because taking the limit
as $h \to 0$ tests inputs $a+h$ on both the left and right; if $f$ were missing
points on either side, the two-sided limit could not even be probed.)
\end{remark}

\begin{remark}[A second, equivalent form]
\label{ch10:def-equiv}
Sometimes it is cleaner to name the moving point directly. Set $x = a + h$, so
$h = x - a$, and notice that $h \to 0$ is the same as $x \to a$. Then
\[
  f'(a) \;=\; \lim_{x \to a} \frac{f(x) - f(a)}{x - a}.
\]
This is the identical definition with a change of variable; use whichever is more
convenient for the algebra in front of you.
\end{remark}

\begin{focusbox}
Read Definition~\ref{ch10:def-derivative} as an \emph{instruction}, not a formula
to plug into. It says: form the difference quotient (a function of $h$), simplify
it until the offending $h$ in the denominator is gone, and \emph{then} see what it
approaches as $h \to 0$. The art of computing derivatives from the definition is
entirely the art of making that $h$ in the denominator cancel.
\end{focusbox}

\section{Computing derivatives from first principles}
\label{ch10:first-principles}

Let us apply the definition directly, before introducing shortcut rules -- just algebra and a
limit. (The differentiation rules of the next chapter all get \emph{proved} this way, so
this is the foundation the rules stand on.)

\begin{example}[A parabola]
\label{ch10:ex-square}
Find the derivative of $f(x) = x^2$ at a general point $x$.
\end{example}

\noindent We build the difference quotient and chase the cancellation:
\begin{align*}
  \frac{f(x + h) - f(x)}{h}
  &= \frac{(x + h)^2 - x^2}{h} \\[2pt]
  &= \frac{x^2 + 2xh + h^2 - x^2}{h}
   = \frac{2xh + h^2}{h} \\[2pt]
  &= \frac{h(2x + h)}{h}
   = 2x + h \qquad (h \neq 0).
\end{align*}
The $h$ cancelled. Now take the limit:
\[
  f'(x) = \lim_{h \to 0} (2x + h) = 2x.
\]
So the derivative of $x^2$ is $2x$. Sanity check against
Figure~\ref{ch10:fig-secants}, which used $\tfrac12 x^2$: at $x = 1$ the slope
there should be $\tfrac12 \cdot 2 \cdot 1 = 1$, and indeed the tangent we drew had
slope $1$. The picture and the algebra agree.

\begin{labbox}
Do this one entirely by hand before reading the solution. Let $f(x) = 3x^2 - 5x
+ 1$. Form $\dfrac{f(x+h) - f(x)}{h}$, expand, cancel the $h$, and take the limit.

\smallskip
\noindent\emph{Solution sketch.} The numerator is
$3(x+h)^2 - 5(x+h) + 1 - (3x^2 - 5x + 1) = 6xh + 3h^2 - 5h$. Divide by $h$ to get
$6x + 3h - 5$, and let $h \to 0$: the derivative is $6x - 5$. Notice the constant
$+1$ contributed nothing -- a constant has zero rate of change, because shifting
the whole graph up and down never tilts it.
\end{labbox}

\begin{example}[A reciprocal -- cancellation by a common denominator]
\label{ch10:ex-recip}
Find the derivative of $f(x) = \dfrac{1}{x}$ for $x \neq 0$.
\end{example}

\noindent Here the cancellation comes from combining fractions, not expanding a
square:
\begin{align*}
  \frac{f(x+h) - f(x)}{h}
  &= \frac{1}{h}\left( \frac{1}{x + h} - \frac{1}{x} \right)
   = \frac{1}{h} \cdot \frac{x - (x + h)}{(x + h)\,x} \\[4pt]
  &= \frac{1}{h} \cdot \frac{-h}{(x + h)\,x}
   = \frac{-1}{(x + h)\,x} \qquad (h \neq 0).
\end{align*}
Again the $h$ vanished, this time absorbed into the $-h$ on top. Now let $h \to 0$:
\[
  f'(x) = \lim_{h \to 0} \frac{-1}{(x + h)\,x} = \frac{-1}{x \cdot x} = -\frac{1}{x^2}.
\]
The minus sign reflects the geometry: $1/x$ slopes \emph{downward}
everywhere it is defined, so its rate of change is negative, and the steeper plunge
near $x = 0$ shows up as the $x^2$ blowing up in the denominator.

\begin{example}[A square root -- cancellation by the conjugate]
\label{ch10:ex-sqrt}
Find the derivative of $f(x) = \sqrt{x}$ for $x > 0$.
\end{example}

\noindent The trick when a root blocks the cancellation is to multiply top and
bottom by the \textbf{conjugate}, $\sqrt{x+h} + \sqrt{x}$ -- the identical pair
of terms but with the minus sign flipped to plus. This exploits the difference of
squares $(A - B)(A + B) = A^2 - B^2$ to strip away the radicals:
\begin{align*}
  \frac{\sqrt{x + h} - \sqrt{x}}{h}
  &= \frac{\sqrt{x + h} - \sqrt{x}}{h}
     \cdot \frac{\sqrt{x + h} + \sqrt{x}}{\sqrt{x + h} + \sqrt{x}} \\[4pt]
  &= \frac{(x + h) - x}{h\bigl(\sqrt{x + h} + \sqrt{x}\bigr)}
   = \frac{h}{h\bigl(\sqrt{x + h} + \sqrt{x}\bigr)} \\[4pt]
  &= \frac{1}{\sqrt{x + h} + \sqrt{x}} \qquad (h \neq 0).
\end{align*}
The lone $h$ on top cancelled the $h$ in the denominator. Now the limit is safe,
because $\sqrt{x+h} \to \sqrt{x}$ as $h \to 0$:
\[
  f'(x) = \lim_{h \to 0} \frac{1}{\sqrt{x + h} + \sqrt{x}}
        = \frac{1}{2\sqrt{x}}.
\]

\begin{whybox}
Why does the conjugate trick work, and why do we keep reaching for these
maneuvers? Every one of them -- expanding $(x+h)^2$, finding a common
denominator, multiplying by the conjugate -- is a way to expose a hidden factor
of $h$ in the numerator so it can cancel the $h$ in the denominator. The
difference quotient is $\frac00$ at $h=0$ precisely \emph{because} the numerator
$f(x+h) - f(x)$ also vanishes there (it is the difference of a quantity with
itself when $h=0$). A quantity that vanishes at $h=0$ has $h$ as a factor, at
least to first order. The algebra simply drags that factor out into the open where
it can be cancelled. Differentiating from first principles is the systematic art
of factoring out the $h$ that \emph{must} be there.
\end{whybox}

\begin{trapbox}
The single most common error is setting $h = 0$ \emph{too early}. If you write
$\dfrac{f(x+h)-f(x)}{h}$ and immediately substitute $h = 0$, you get
$\dfrac{0}{0}$ and learn nothing -- worse, you may convince yourself the
derivative ``doesn't exist'' when it is perfectly fine. The rule is simple:
\emph{simplify until the denominator no longer contains $h$, and only then evaluate the limit as
$h \to 0$}. Cancellation first, limit second. Never the other way.
\end{trapbox}

\begin{example}[A constant and a line -- the easy extremes]
\label{ch10:ex-const-line}
If $f(x) = c$ is constant, then $f(x+h) - f(x) = c - c = 0$, so the difference
quotient is $0/h = 0$ for every $h \neq 0$, and $f'(x) = 0$. A flat graph has zero
slope; a quantity that never changes has zero rate of change. If $f(x) = mx + b$
is a line, then
\[
  \frac{m(x+h)+b - (mx+b)}{h} = \frac{mh}{h} = m,
\]
so $f'(x) = m$ for every $x$. The derivative of a line is its slope, constant
everywhere -- which had better be true, since the tangent to a straight line is
the line itself.
\end{example}

\section{The derivative as a function}
\label{ch10:deriv-as-function}

Look back at the examples. We started out asking for the slope at one point $a$,
but the algebra naturally handed us a \emph{formula in $x$}: the slope of $x^2$ is
$2x$, valid at every $x$. This is an essential shift in perspective. The derivative is not
just a number attached to one point -- it is a \emph{new function} that reports
the slope at every point where the original is differentiable.

\begin{definition}[The derivative function]
\label{ch10:def-deriv-function}
Given $f$, its \textbf{derivative function} $f'$ is defined by
\[
  f'(x) \;=\; \lim_{h \to 0} \frac{f(x + h) - f(x)}{h},
\]
on the set of all $x$ for which the limit exists. The process of producing $f'$
from $f$ is called \textbf{differentiation}.
\end{definition}

There is genuine information in graphing $f$ and $f'$ together. Where $f$ rises,
$f'$ is positive; where $f$ falls, $f'$ is negative; where $f$ momentarily levels
off at a peak or valley, $f'$ crosses zero. Figure~\ref{ch10:fig-f-and-fprime}
shows $f(x) = x^2$ over its derivative $f'(x) = 2x$: the parabola's slope is
negative on the left, zero at the bottom, positive on the right, and the straight
line $2x$ records exactly that.

\begin{figure}[htbp]
\centering
\begin{tikzpicture}
\begin{axis}[
  width=12cm, height=7cm,
  title style={font=\small},
  title={Function versus derivative: slope values plotted as heights},
  axis lines=middle,
  xlabel={$x$}, ylabel={},
  xmin=-2.4, xmax=2.4, ymin=-4.5, ymax=4.8,
  xtick={-2,-1,1,2}, ytick={-4,-2,2,4},
  domain=-2.2:2.2, samples=100,
  legend pos=south east,
  legend cell align=left,
]
\addplot[curveblue, very thick] {x^2};
\addlegendentry{$f(x)=x^2$}
\addplot[tangentred, very thick, dashed] {2*x};
\addlegendentry{$f'(x)=2x$}
\addplot[gray, dotted, thick] coordinates {(-2.4,0)(2.4,0)};
\end{axis}
\end{tikzpicture}
\caption{A function and its derivative. The slope of $f(x)=x^2$ at each $x$ is
read off the line $f'(x)=2x$: negative for $x<0$ (the parabola falls), zero at
$x=0$ (the bottom of the bowl), positive for $x>0$ (the parabola rises). The
height of the dashed line \emph{is} the steepness of the solid curve.}
\label{ch10:fig-f-and-fprime}
\end{figure}

\begin{connectionbox}
This is the seed of an entire subject -- the qualitative analysis of curves
(Chapter 16). The sign of $f'$ tells you whether $f$ is increasing or decreasing;
the zeros of $f'$ locate the candidate peaks and valleys. We will eventually
differentiate $f'$ \emph{again} to get $f''$, the rate at which the rate is
changing -- which is acceleration in the falling-ball story and concavity in the
geometric one. The derivative being a function, and not merely a single number, is what makes
all of that possible.
\end{connectionbox}

\section{The third meaning: best linear approximation}
\label{ch10:best-linear}

We have met two meanings of the derivative -- slope of the tangent and
instantaneous rate of change. The third is subtler, and it carries most of the
weight in the rest of calculus. It says: \emph{near $a$, the function $f$ looks
like a straight line, and the derivative is the slope of that line.}

Here is the idea. Unwind the definition. To say
\[
  \lim_{h \to 0} \frac{f(a + h) - f(a)}{h} = f'(a)
\]
is to say that for small $h$, the difference quotient is approximately $f'(a)$:
\[
  \frac{f(a + h) - f(a)}{h} \approx f'(a),
  \qquad\text{hence}\qquad
  f(a + h) \approx f(a) + f'(a)\,h.
\]
Writing $x = a + h$ (so $h = x - a$), this is the \textbf{tangent line
approximation}
\[
  \boxed{\,f(x) \approx f(a) + f'(a)\,(x - a)\,}
\]
valid for $x$ near $a$. The right-hand side is precisely the equation of the
tangent line at $a$: it passes through $(a, f(a))$ and has slope $f'(a)$. So the
tangent line is not merely a line that grazes the curve -- it is the \emph{best
straight-line stand-in} for $f$ near $a$. Among all lines through $(a, f(a))$, the
tangent is the one whose error shrinks faster than $h$ itself as $h \to 0$.

\begin{intuitionbox}
Zoom in. Take any smooth curve, fix a point on it, and magnify the picture around
that point -- double the magnification, then double again, and again. The curve
flattens. At high enough zoom it is visually indistinguishable from a straight
line, and that line is the tangent. This is \textbf{local linearity}, and it is
the real reason calculus works. Curves are hard; lines are easy. The derivative is
the dictionary that translates the hard local behavior of a curve into the slope
of a line. Almost every technique you will learn -- linear approximation,
Newton's method, and the methods of differential calculus -- is just this zoom,
exploited. Figure~\ref{ch10:fig-zoom} performs the zoom on a genuine curve.
\end{intuitionbox}

\begin{figure}[htbp]
\centering
\begin{tikzpicture}
% Panel 1: wide view
\begin{axis}[
  width=4.7cm, height=4.7cm, name=zA,
  title style={font=\small},
  axis lines=middle, title={wide view},
  xmin=-1.6, xmax=1.6, ymin=-0.4, ymax=2.6,
  xtick=\empty, ytick=\empty, domain=-1.5:1.5, samples=80,
]
\addplot[curveblue, very thick] {1 + 0.5*x + 0.5*x^2};
\addplot[tangentred, very thick, dashed] {1 + 0.5*x};
\addplot[only marks, mark=*, mark size=1.3pt, black] coordinates {(0,1)};
\end{axis}
% Panel 2: zoom x10
\begin{axis}[
  width=4.7cm, height=4.7cm, name=zB,
  at={($(zA.east)+(0.8cm,0)$)}, anchor=west,
  title style={font=\small},
  axis lines=middle, title={zoom $\times 10$},
  xmin=-0.16, xmax=0.16, ymin=0.84, ymax=1.16,
  xtick=\empty, ytick=\empty, domain=-0.15:0.15, samples=80,
]
\addplot[curveblue, very thick] {1 + 0.5*x + 0.5*x^2};
\addplot[tangentred, very thick, dashed] {1 + 0.5*x};
\addplot[only marks, mark=*, mark size=1.3pt, black] coordinates {(0,1)};
\end{axis}
% Panel 3: zoom x100
\begin{axis}[
  width=4.7cm, height=4.7cm,
  at={($(zB.east)+(0.8cm,0)$)}, anchor=west,
  title style={font=\small},
  axis lines=middle, title={zoom $\times 100$},
  xmin=-0.016, xmax=0.016, ymin=0.984, ymax=1.016,
  xtick=\empty, ytick=\empty, domain=-0.015:0.015, samples=80,
]
\addplot[curveblue, very thick] {1 + 0.5*x + 0.5*x^2};
\addplot[tangentred, very thick, dashed] {1 + 0.5*x};
\addplot[only marks, mark=*, mark size=1.3pt, black] coordinates {(0,1)};
\end{axis}
\end{tikzpicture}
\caption{Local linearity, made literal. The same curve
$f(x)=1+\tfrac12 x+\tfrac12 x^2$ (solid) and its tangent at $(0,1)$ (dashed),
viewed at three magnifications about that point. As we zoom, the curve becomes
indistinguishable from its tangent: by $\times 100$ the two are visually one
line. The derivative $f'(0)=\tfrac12$ is the slope of the line the curve turns
into under the microscope.}
\label{ch10:fig-zoom}
\end{figure}

\begin{figure}[htbp]
\centering
\begin{tikzpicture}
\begin{axis}[
  width=12cm, height=6.5cm,
  title style={font=\small},
  title={Best linear approximation: $\sqrt{x}$ and its tangent at $x=1$},
  axis lines=middle,
  xlabel={$x$}, ylabel={$y$},
  xmin=0.2, xmax=1.85, ymin=0, ymax=1.7,
  xtick={1}, ytick={1},
  domain=0.25:1.8, samples=120,
  legend pos=north west, legend cell align=left,
]
% f(x) = sqrt(x), tangent at x=1: y = 1 + (1/2)(x-1)
\addplot[curveblue, very thick] {sqrt(x)};
\addlegendentry{$f(x)=\sqrt{x}$}
\addplot[tangentred, very thick, dashed] {1 + 0.5*(x-1)};
\addlegendentry{tangent at $x=1$}
\addplot[only marks, mark=*, mark size=2pt, black] coordinates {(1,1)};
\node[below right] at (axis cs:1,1) {$(1,1)$};
\end{axis}
\end{tikzpicture}
\caption{The tangent line as best linear approximation. Near $x=1$ the curve
$\sqrt{x}$ and its tangent $y = 1 + \tfrac12(x-1)$ are nearly on top of each
other; they separate only as you move away from the point of contact. The
approximation $\sqrt{x} \approx 1 + \tfrac12(x-1)$ gives, e.g.,
$\sqrt{1.04} \approx 1.02$.}
\label{ch10:fig-linear}
\end{figure}

\begin{example}[Approximating a square root in your head]
\label{ch10:ex-approx}
From Example~\ref{ch10:ex-sqrt}, $f(x) = \sqrt{x}$ has $f'(x) = \tfrac{1}{2\sqrt
x}$. Near $a = 1$ we have $f(1) = 1$ and $f'(1) = \tfrac12$, so
\[
  \sqrt{x} \approx 1 + \tfrac12 (x - 1).
\]
Then $\sqrt{1.04} \approx 1 + \tfrac12(0.04) = 1.02$. The true value is
$1.0198\ldots$, so we are right to three decimal places with mental arithmetic and
no calculator -- the payoff of replacing a curve by its tangent line. The chapter
on linear approximation and differentials develops this systematically.
\end{example}

\begin{formalbox}
Here is the precise statement that justifies calling the tangent the ``best''
linear approximation. Define the error of the tangent line at $a$ as
\[
  E(h) = f(a + h) - \bigl[\, f(a) + f'(a)\,h \,\bigr].
\]
Then differentiability at $a$ is \emph{equivalent} to the statement
\[
  \lim_{h \to 0} \frac{E(h)}{h} = 0,
\]
that is, the error vanishes faster than $h$ itself. Any \emph{other} line through
$(a,f(a))$, with slope $m \neq f'(a)$, has error $\bigl(f'(a) - m\bigr) h + E(h)$,
whose ratio to $h$ tends to the nonzero number $f'(a) - m$. Only the tangent's
error is negligible compared to $h$. This ``error smaller than first order'' is
the rigorous meaning of ``best,'' and it is the doorway to Taylor's theorem later
in the book.
\end{formalbox}

\section{Notation: Leibniz, Lagrange, and what $\dv{y}{x}$ really means}
\label{ch10:notation}

Calculus was invented twice, more or less simultaneously, by two people who could
hardly have been more different, and we inherited \emph{both} of their notations.
It pays to know each, because each makes certain truths obvious and others
awkward.

\begin{historybox}
The history of calculus involves two independent inventors, a prolonged priority
dispute, foundational philosophical challenges, and eventual formalization.

In the 1660s in England, \textbf{Isaac Newton} was driven by physics: how do
planets orbit the sun, and how does a falling apple accelerate? He imagined
quantities ``flowing'' through continuous time. He called a changing quantity a
\textbf{fluent} (a quantity that flows), and its instantaneous rate of change a
\textbf{fluxion} (its flowing speed), writing $\dot{x}$ for the fluxion of $x$
-- a dot notation still prized in physics today. A decade later in Germany,
\textbf{Gottfried Wilhelm Leibniz}, working independently from 1675 to 1684,
approached the problem as a philosopher of language. He wanted a notation so
transparent that algebraic calculations would automatically mirror geometric
truths. He wrote the derivative as $\dv{y}{x}$ -- a ratio of two
\textbf{infinitesimals} (idealized quantities smaller than any positive real
number, yet strictly not zero).

Both thinkers were stuck on the identical foundational hurdle: calculating an
instantaneous rate requires dividing a change in output by an elapsed time step,
yet at a single frozen moment both changes are zero, producing the meaningless
expression $\frac{0}{0}$. Both relied on Cartesian coordinate geometry and
clever algebraic tricks, but neither possessed a rigorous theory of limits.
Instead, they divided by tiny infinitesimal increments and then discarded the
leftover pieces at the end.

Two major controversies erupted. First was the \textbf{priority dispute}
-- a bitter public battle over who discovered calculus first. Newton had developed
his ideas first in private manuscripts, but Leibniz was the first to publish in
1684. The Royal Society of London and Continental mathematicians exchanged sharp
accusations of plagiarism. The dispute alienated British mathematicians from the
Continent for more than a century; loyalty to Newton's fluxion notation
caused British mathematics to lag, while Continental scholars advanced
using Leibniz's differential notation.

The second crisis concerned the logical foundations. In 1734, the philosopher Bishop
\textbf{George Berkeley} published a critique titled \emph{The Analyst}.
Berkeley exposed a central gap in early derivative calculations:
in simplifying a difference quotient, one first assumes that the step $h$ is
nonzero in order to divide by it, and then sets $h = 0$ to drop the remaining
terms, relying on two incompatible assumptions in a single argument.

Despite Berkeley's critique, calculus provided the mathematical foundation for
celestial mechanics, planetary orbits, ocean tides, and differential equations
governing heat, wave motion, and fluid flow. The rigorous resolution arrived in the
nineteenth century through the work of \textbf{Augustin-Louis Cauchy} and
\textbf{Karl Weierstrass}, who placed calculus on the foundation of limits
and the $\eps$--$\delta$ definition (Chapters 5--7). In their system, $h$ is never
set to zero; we simply evaluate the limit as $h$ approaches zero.

In the 1960s, mathematical logician \textbf{Abraham Robinson} developed
\textbf{nonstandard analysis} -- a branch of mathematical logic that uses
formal structures to construct the \emph{hyperreal numbers}, a number system
containing genuine, logically sound infinitesimals. Robinson proved that
infinitesimal arguments can be made completely rigorous and
reproduce all results of standard calculus.
\end{historybox}

\textbf{Lagrange's prime notation.} Joseph-Louis Lagrange later introduced the
clean $f'(x)$ we have been using -- ``$f$ prime of $x$.'' It is compact, it clearly
emphasizes the derivative-as-a-function ($f'$ is a function just like $f$), and it
is well suited for stating theorems. Its weakness: it hides \emph{what} you are
differentiating with respect to. For $f'$ that is fine; for a quantity depending
on several variables it is ambiguous.

\textbf{Leibniz's ratio notation.} If we name the output $y = f(x)$, then Leibniz
writes
\[
  \dv{y}{x} \quad\text{or}\quad \frac{\dd}{\dd x}f(x) \quad\text{or}\quad
  \ddx{f(x)},
\]
and the value at a specific point $a$ is written
$\left.\dv{y}{x}\right|_{x=a}$. The symbol $\dfrac{\dd}{\dd x}$ is an
\emph{operator}: ``take the derivative with respect to $x$ of whatever follows.''
Leibniz notation is wonderful because it \emph{remembers the variables}, it makes
the chain rule and substitution look like fraction arithmetic (as we will see),
and it makes units obvious -- $\dv{s}{t}$ visibly carries units of distance over
time, i.e.\ a speed.

\begin{happeningbox}
Here is the question that trips up nearly every student, so let us answer it
squarely. \emph{Is $\dv{y}{x}$ a fraction?}

In the way we have built it, \textbf{no}. The symbol $\dv{y}{x}$ is a single,
indivisible name for one number: the limit $\displaystyle\lim_{h\to 0}\frac{f(x+h)
-f(x)}{h}$. There is no genuine number called $\dd y$ sitting on top of a genuine
number $\dd x$ that you divided. Leibniz himself thought of $\dd x$ and $\dd y$ as
``infinitesimals'' -- quantities smaller than any positive number yet not zero --
and treated $\dv{y}{x}$ as their literal quotient. For two centuries no one could
say rigorously what an infinitesimal \emph{was}, and the foundations were, frankly,
a mess (see the History box above).

So why does the fraction notation work so flawlessly? Because $\dv{y}{x}$ is the
\emph{limit} of ordinary fractions $\frac{\Delta y}{\Delta x} = \frac{f(x+h)-f(x)}
{h}$, and many of the algebraic moves you would make on the fraction $\frac{\Delta
y}{\Delta x}$ \emph{survive the limit}. The chain rule $\dv{y}{x} = \dv{y}{u}\cdot
\dv{u}{x}$ looks like cancelling $\dd u$, and it is \emph{true} -- but it is true
because we will \emph{prove} it from the definition, not because $\dd u$ literally
cancels. The notation is a brilliantly suggestive mnemonic for theorems that are
established by limits. Use the fraction intuition to \emph{guess}; lean on the
limit definition to \emph{know}.
\end{happeningbox}

\section{Differentiability versus continuity}
\label{ch10:diff-vs-cont}

A function with a derivative at $a$ is rather special, and it helps to know exactly
\emph{how} special. There is a clean one-way implication.

\begin{theorem}[Differentiable $\Rightarrow$ continuous]
\label{ch10:thm-diff-implies-cont}
If $f$ is differentiable at $a$, then $f$ is continuous at $a$.
\end{theorem}

\begin{intuitionbox}
Before the proof, the picture. If $f$ has a tangent line at $a$, then near $a$ the
graph hugs that line -- and a line has no gaps or jumps. So the graph cannot
suddenly leap at $a$; it must arrive at $f(a)$ smoothly. Having a well-defined
slope is a \emph{stronger} demand than merely having no jump. Differentiability
implies continuity, but, as we will see, continuity does not imply
differentiability.
\end{intuitionbox}

\begin{proof}
We must show $\displaystyle\lim_{x \to a} f(x) = f(a)$, equivalently
$\displaystyle\lim_{x \to a}\bigl[f(x) - f(a)\bigr] = 0$. The trick is to
manufacture the difference quotient by multiplying and dividing by $x - a$ (legal
for $x \neq a$, which is all the limit ever sees):
\[
  f(x) - f(a) = \frac{f(x) - f(a)}{x - a}\,\cdot\,(x - a).
\]
Now take the limit as $x \to a$. The first factor tends to $f'(a)$ by the
definition of the derivative (Remark~\ref{ch10:def-equiv}), which exists by
hypothesis; the second factor tends to $0$. By the product law for limits,
\[
  \lim_{x \to a}\bigl[f(x) - f(a)\bigr]
  = f'(a) \cdot 0 = 0.
\]
Hence $\lim_{x \to a} f(x) = f(a)$, which is exactly continuity at $a$.
\end{proof}

The converse is false, and its failure is so instructive that we devote the rest
of the section to the three classic ways a continuous function can lack a
derivative. In every case the function has no jump -- it is continuous -- yet the
limit defining $f'(a)$ does not exist. Watch \emph{why} the limit fails each time;
it is always the secant slopes refusing to settle on one number.

\subsection*{Failure 1: a corner -- the absolute value}

The cleanest example is $f(x) = \abs{x}$ at $a = 0$. It is continuous everywhere
(no jump at $0$; $\abs{x} \to 0 = \abs{0}$). But form the difference quotient at
$0$:
\[
  \frac{f(0 + h) - f(0)}{h} = \frac{\abs{h} - 0}{h} = \frac{\abs{h}}{h}
  = \begin{cases} \;\;1 & h > 0, \\[2pt] -1 & h < 0. \end{cases}
\]
As $h \to 0^+$ the quotient is $+1$; as $h \to 0^-$ it is $-1$. The one-sided
limits disagree, so the two-sided limit does not exist, and $f'(0)$ is undefined.
Geometrically the graph has a sharp \textbf{corner} at the origin: the slope is
$-1$ coming from the left and $+1$ going to the right, and there is no single
tangent line. The secants from the right insist the slope is $1$; the secants from
the left insist it is $-1$; they never agree.

\subsection*{Failure 2: a cusp -- the cube root, sharpened}

Consider $f(x) = x^{2/3}$ at $a = 0$. It is continuous ($0^{2/3} = 0$), and unlike
the corner its two sides do not meet at an angle -- they meet at an infinitely
sharp \textbf{cusp} -- a pointed tip where the curve reverses direction and both
branches turn vertical. The difference quotient is
\[
  \frac{(0+h)^{2/3} - 0}{h} = \frac{h^{2/3}}{h} = h^{-1/3} = \frac{1}{h^{1/3}},
\]
and as $h \to 0$ this grows without bound (toward $+\infty$ from the right, toward
$-\infty$ from the left). The limit does not exist -- not even as one-sided finite
limits -- so $f'(0)$ does not exist. The slope is trying to be vertical on both
sides at once.

\subsection*{Failure 3: a vertical tangent -- the cube root}

Now $f(x) = x^{1/3}$ at $a = 0$, again continuous. The difference quotient is
\[
  \frac{(0+h)^{1/3} - 0}{h} = \frac{h^{1/3}}{h} = h^{-2/3} = \frac{1}{h^{2/3}}
  \;\To\; +\infty \quad \text{as } h \to 0.
\]
This time both sides agree -- the quotient marches to $+\infty$ from either
direction -- so the curve genuinely has a single limiting tangent direction, but
that direction is \emph{vertical}. A vertical line has no (finite) slope, so by
our definition $f'(0)$ does not exist as a number. We say the graph has a
\textbf{vertical tangent} at the origin. It is differentiability's near miss: there
is a tangent line, but its slope is infinite.

\begin{figure}[htbp]
\centering
\begin{tikzpicture}
\begin{axis}[
  width=4.7cm, height=4.7cm, name=plotA,
  title style={font=\small},
  axis lines=middle, title={corner: $\abs{x}$},
  xmin=-1.4, xmax=1.4, ymin=-0.3, ymax=1.4,
  xtick=\empty, ytick=\empty, domain=-1.3:1.3, samples=80,
]
\addplot[curveblue, very thick] {abs(x)};
\addplot[only marks, mark=*, mark size=1.5pt, tangentred] coordinates {(0,0)};
\end{axis}
\begin{axis}[
  width=4.7cm, height=4.7cm, name=plotB,
  at={($(plotA.east)+(0.8cm,0)$)}, anchor=west,
  title style={font=\small},
  axis lines=middle, title={cusp: $x^{2/3}$},
  xmin=-1.4, xmax=1.4, ymin=-0.3, ymax=1.4,
  xtick=\empty, ytick=\empty, domain=-1.3:1.3, samples=120,
]
\addplot[curveblue, very thick] {abs(x)^(2/3)};
\addplot[only marks, mark=*, mark size=1.5pt, tangentred] coordinates {(0,0)};
\end{axis}
\begin{axis}[
  width=4.7cm, height=4.7cm,
  at={($(plotB.east)+(0.8cm,0)$)}, anchor=west,
  title style={font=\small},
  axis lines=middle, title={vertical: $x^{1/3}$},
  xmin=-1.4, xmax=1.4, ymin=-1.4, ymax=1.4,
  xtick=\empty, ytick=\empty, domain=-1.3:1.3, samples=120,
]
\addplot[curveblue, very thick] {sign(x)*abs(x)^(1/3)};
\addplot[only marks, mark=*, mark size=1.5pt, tangentred] coordinates {(0,0)};
\end{axis}
\end{tikzpicture}
\caption{Three ways to be continuous but not differentiable at the origin. A
\emph{corner} ($\abs{x}$): left and right slopes disagree. A \emph{cusp}
($x^{2/3}$): the secant slopes run off to $+\infty$ on one side and $-\infty$ on
the other. A \emph{vertical tangent} ($x^{1/3}$): a genuine tangent exists, but it
is vertical, so its slope is infinite.}
\label{ch10:fig-failures}
\end{figure}

\begin{walkbox}
\textbf{Pause and think:} Consider $f(x) = \abs{x - 3}$. At what point does the
corner occur, and what are the left-hand and right-hand slopes there? Without
writing any algebra, what are $f'(5)$ and $f'(1)$? (The corner is at $x=3$, where
slopes abruptly jump from $-1$ to $+1$. For $x > 3$, the function is the straight
line $y = x - 3$, so $f'(5) = 1$. For $x < 3$, it is the line $y = -(x - 3)$, so
$f'(1) = -1$. The derivative exists everywhere except at the single corner point
$x=3$.)
\end{walkbox}

\begin{trapbox}
Do not over-learn Theorem~\ref{ch10:thm-diff-implies-cont}. It runs
\emph{one way only}: differentiable $\Rightarrow$ continuous. The reverse is false,
as $\abs{x}$ shows -- continuous everywhere, not differentiable at $0$. So
``continuous'' is necessary but not sufficient for ``differentiable.'' A handy
slogan: \emph{differentiable functions are continuous and smooth; continuous
functions are merely unbroken, and may still be jagged.}
\end{trapbox}

\begin{aestheticbox}
How bad can the gap between continuous and differentiable get? In 1872 Weierstrass
exhibited a function that is \emph{continuous everywhere yet differentiable nowhere}
-- an unbroken curve that has a corner at every single point, so jagged that no point
on it admits a tangent line. When it was published it surprised many mathematicians;
Charles Hermite referred to such functions as an unnatural pathology. Today they
are fundamental to fractal geometry and the mathematics of Brownian motion -- the
erratic path of a microscopic particle suspended in fluid, which almost surely traces
an unbroken path that fails to have a well-defined slope anywhere. Continuity and
differentiability are genuinely different properties.
\end{aestheticbox}

\section{Rates of change everywhere}
\label{ch10:rates-everywhere}

We opened with a falling ball because motion makes the derivative tangible, but
the derivative is the universal language of \emph{any} quantity changing with
respect to \emph{any} other. The pattern is always the same: a small change in the
input causes a small change in the output, and the derivative is the limiting
ratio.

\begin{applicationbox}
The derivative goes by a different name in every field, but it is always
$\displaystyle\lim_{h\to 0}\frac{f(x+h)-f(x)}{h}$.
\begin{itemize}
  \item \textbf{Physics.} If $s(t)$ is position, $s'(t)$ is velocity, and
  $s''(t)$ is acceleration. Newton's second law $F = ma$ is a statement about the
  second derivative of position.
  \item \textbf{Biology.} If $P(t)$ is a population, $P'(t)$ is the growth rate
  (births minus deaths per unit time). The model $P'(t) = kP(t)$ -- growth rate
  proportional to current size -- gives exponential growth, and is your first
  differential equation.
  \item \textbf{Economics.} If $C(q)$ is the cost of producing $q$ items, the
  derivative $C'(q)$ is the \textbf{marginal cost} -- the additional expense
  incurred by producing one single additional item. ``Marginal'' in economics
  almost always means ``derivative.''
  \item \textbf{Chemistry.} If $[A](t)$ is a concentration, the reaction rate is
  $-\dv{[A]}{t}$ -- the speed at which reactant is consumed.
  \item \textbf{Machine learning.} Training a neural network means minimizing a loss
  function $L(\mathbf{w})$ over parameters $\mathbf{w}$; gradient descent steps
  downhill along the negative \textbf{gradient} (the multi-dimensional vector
  pointing along the direction of steepest increase). Modern AI models adjust
  billions of weights using \textbf{backpropagation} -- an algorithm that
  evaluates the chain rule in reverse to compute derivatives of the output error
  with respect to every internal parameter.
\end{itemize}
The instantaneous rate of change is not a physics trick. It is the way to talk
about change itself, in any domain where one quantity depends on another.
\end{applicationbox}

\begin{walkbox}
Take a walk and find the derivatives hiding around you. The speedometer in a car
literally displays $\dv{s}{t}$. The slope you feel in your legs on a hill is the
derivative of elevation with respect to horizontal distance. The rate your coffee
cools is the derivative of its temperature with respect to time (and, as Newton
noticed, that rate is proportional to how much hotter the coffee is than the room
-- another differential equation). Ask of each: what is changing, with respect to
what, and is the rate positive, negative, or zero right now?
\end{walkbox}

\section{What to carry forward}
\label{ch10:summary}

\begin{recoverybox}
A summary of the chapter's core ideas:
\begin{itemize}
  \item The \textbf{derivative} of $f$ at $a$ is
  $f'(a) = \displaystyle\lim_{h\to 0}\frac{f(a+h)-f(a)}{h}$, when this limit
  exists. It is the limit of difference quotients -- of average rates -- as the
  interval shrinks to nothing.
  \item It has \textbf{three faces}, all one number: the \emph{slope of the
  tangent line}, the \emph{instantaneous rate of change}, and the \emph{slope of
  the best linear approximation} $f(x)\approx f(a)+f'(a)(x-a)$.
  \item Compute it from first principles by \textbf{cancelling the $h$ first}
  (expand, common-denominator, or conjugate) and \textbf{taking the limit second}.
  Never set $h=0$ before the cancellation.
  \item The derivative is itself a \textbf{function}, $f'$, recording the slope at
  every point.
  \item \textbf{Notation}: Lagrange's $f'(x)$ and Leibniz's $\dv{y}{x}$ name the
  same thing. $\dv{y}{x}$ is \emph{not} literally a fraction -- it is the limit of
  ordinary fractions, which is why fraction-style manipulations keep working once we
  prove them.
  \item \textbf{Differentiable $\Rightarrow$ continuous}, but not conversely:
  corners ($\abs{x}$), cusps ($x^{2/3}$), and vertical tangents ($x^{1/3}$) are
  continuous yet have no derivative, because the secant slopes refuse to settle on
  one finite number.
\end{itemize}
Next chapter we stop computing every derivative from scratch and build the
\emph{rules} -- and each rule we prove will rest squarely on the limit definition
you now own.
\end{recoverybox}

\section{Exercises}
\label{ch10:exercises}

\begin{exercise}[Mental checks on sight]
\label{ch10:ex-mental-checks}
Answer each of these in your head in ten seconds without pencil or paper:
\begin{enumerate}[label=(\alph*)]
  \item What is the derivative of the constant function $f(x) = 42$? What is the slope of its graph?
  \item What is the derivative of the straight line $f(x) = -7x + 3$?
  \item If $s(t)$ measures a car's distance in miles and $t$ is in hours, what are the units of $\dv{s}{t}$?
  \item If the tangent line to $y = f(x)$ at $x = 5$ is horizontal, what is $f'(5)$?
  \item Given $f(2) = 10$ and $f'(2) = 3$, estimate $f(2.04)$ using linear approximation.
\end{enumerate}
\end{exercise}

\begin{exercise}[Warm-ups from the definition]
\label{ch10:ex-warmups}
Using only Definition~\ref{ch10:def-derivative} (cancel the $h$, then take the
limit), compute $f'(x)$ for each:
\begin{enumerate}[label=(\alph*)]
  \item $f(x) = 5x - 7$.
  \item $f(x) = x^2 - 3x$.
  \item $f(x) = x^3$. \emph{(Hint: $(x+h)^3 = x^3 + 3x^2 h + 3x h^2 + h^3$.)}
  \item $f(x) = \dfrac{1}{x + 1}$.
\end{enumerate}
\end{exercise}

\begin{exercise}[A root, by the conjugate]
\label{ch10:ex-root}
Compute the derivative of $f(x) = \sqrt{2x + 1}$ from first principles, using the
conjugate trick. At what value of $x$ is the slope equal to $1$?
\end{exercise}

\begin{exercise}[Instantaneous speed]
\label{ch10:ex-speed}
A particle moves so that its position at time $t$ is $s(t) = t^2 - 4t + 5$ (in
meters, $t$ in seconds). (a) Find the average velocity over $[1, 1+h]$ and
simplify. (b) Take the limit to find the instantaneous velocity at $t = 1$.
(c) At what time is the particle instantaneously at rest (velocity zero)? What is
it doing at that instant -- where on its path is it?
\end{exercise}

\begin{exercise}[Tangent line and approximation]
\label{ch10:ex-tangent}
Let $f(x) = x^2$. (a) Write the equation of the tangent line at $a = 3$. (b) Use
it to approximate $(3.02)^2$ without squaring, and compare with the exact value.
(c) Is your approximation an over- or under-estimate? Explain using the shape
(concavity) of the parabola.
\end{exercise}

\begin{exercise}[Reading a graph]
\label{ch10:ex-reading}
A function $f$ is increasing for $x < 2$, has a peak at $x = 2$, and decreases for
$x > 2$. Sketch a plausible graph of the \emph{derivative} $f'$. Where is $f'$
positive, where negative, and what happens at $x = 2$? (No formula is needed --
reason from signs of slopes.)
\end{exercise}

\begin{exercise}[Where the derivative fails]
\label{ch10:ex-fails}
For each function, decide whether it is differentiable at $a = 0$. If not, name
the type of failure (corner, cusp, vertical tangent, or discontinuity) and justify
with one-sided difference quotients.
\begin{enumerate}[label=(\alph*)]
  \item $f(x) = \abs{x - 0} = \abs{x}$ -- restated for contrast with (b).
  \item $f(x) = x\,\abs{x}$. \emph{(Compute the left and right difference
  quotients at $0$; do they agree?)}
  \item $f(x) = x^{1/3}$.
  \item $f(x) = \begin{cases} x & x \ge 0,\\ x - 1 & x < 0.\end{cases}$
\end{enumerate}
\end{exercise}

\begin{exercise}[Continuous but not differentiable -- by hand]
\label{ch10:ex-cont-not-diff}
Let $f(x) = \abs{x}$. (a) Show directly from the limit definition that $f$ is
\emph{continuous} at $0$. (b) Show directly that the two one-sided limits of the
difference quotient at $0$ are $+1$ and $-1$, so $f'(0)$ does not exist. (c) In one
sentence, explain why this does not contradict
Theorem~\ref{ch10:thm-diff-implies-cont}.
\end{exercise}

\begin{exercise}[Prove a special case]
\label{ch10:ex-prove-power}
Using the binomial expansion of $(x+h)^n$, prove from
Definition~\ref{ch10:def-derivative} that for any positive integer $n$,
\[
  \ddx{x^n} = n x^{n-1}.
\]
\emph{Hint: $(x+h)^n = x^n + nx^{n-1}h + (\text{terms with } h^2 \text{ or
higher})$. After cancelling $h$, every surviving term except the first still
contains at least one factor of $h$.}
\end{exercise}

\begin{exercise}[The other definition]
\label{ch10:ex-other-def}
Use the equivalent form $f'(a) = \displaystyle\lim_{x\to a}\dfrac{f(x)-f(a)}{x-a}$
(Remark~\ref{ch10:def-equiv}) to compute $f'(a)$ for $f(x) = x^2$ at a general
$a$, by factoring $x^2 - a^2 = (x-a)(x+a)$. Confirm you get $2a$, matching
Example~\ref{ch10:ex-square}.
\end{exercise}

\begin{exercise}[Conceptual -- is the fraction real?]
\label{ch10:ex-fraction}
A classmate insists that $\dv{y}{x}$ \emph{is} a fraction because the chain rule
``cancels the $\dd u$.'' Write a short paragraph explaining, using the
$\frac00$-and-limit idea from this chapter, in what precise sense $\dv{y}{x}$ is
\emph{not} a fraction, and why fraction-style manipulations nonetheless give
correct answers.
\end{exercise}

\begin{exercise}[A subtle differentiability question]
\label{ch10:ex-subtle}
Consider
\[
  f(x) = \begin{cases} x^2 \sin\!\left(\dfrac{1}{x}\right) & x \neq 0,\\[4pt]
  0 & x = 0.\end{cases}
\]
(a) Using $\abs{\sin(1/x)} \le 1$ and the squeeze theorem, show the difference
quotient $\dfrac{f(h)-f(0)}{h} = h\sin(1/h)$ tends to $0$, so $f'(0) = 0$ exists.
(b) This function is differentiable at $0$ even though it oscillates infinitely
often near $0$. What does this tell you about how much ``wiggling'' is compatible
with having a derivative? \emph{(This classic example shows how the factor $x^2$
bounds the oscillation to yield a well-defined limit.)}
\end{exercise}
